\chapter{Lucas's ten-player game without an NM solution}
\label{sec:p18-lucas-no-solution}

The characteristic function in \cref{def:p18-lucas-game} has a
nonempty core, yet has no von Neumann--Morgenstern stable set of
imputations. This chapter gives the complete regional exclusion
argument of Lucas's detailed report
\cite[secs.~4--8, pp.~6--18]{Lucas1968NoSolutionProof}.
The player set remains the actual ten-player set. Its imputation
space is a continuum; replacing it by a finite selection of
imputations would change the problem.

\begin{hbdefinition}{def:p18-lucas-regional-notation}{Domination and cyclic coordinates}
Throughout this chapter abbreviate
\[
 A=A_\ell,\qquad B=B_\ell,\qquad C={\cal C}_\ell,
 \qquad H=\{1,3,5\}.
\]
The sets, vertices \(c^0,c^2,c^4,c^6,c^8,c^{10}\), and the entire
function \(\ell\) are those of \cref{def:p18-lucas-game}.
In particular,
\[
 B=\{x\in[0,1]^{10}:x_{2a-1}+x_{2a}=1\ (1\leq a\leq5)\},
 \qquad C=\{x\in B:x(O)\geq4\}.
\]
For \(x,y\in A\), write \(y\succ_S x\) when \(S\ne\varnothing\),
\[
 y_h>x_h\quad(h\in S),\qquad y(S)\leq\ell(S);
\]
write \(y\succ x\) when such an \(S\) exists. Thus coalition
feasibility is weak, while every member's improvement is strict.
For \(X\subseteq A\), set
\[
 \operatorname{Dom}X=\{x\in A:\exists y\in X\ (y\succ x)\}.
\]
A solution relative to \(X\) is a subset \(K\subseteq X\) with no
domination between its members and
\(X\setminus K\subseteq\operatorname{Dom}K\).

Let \(\rho(1)=3,\rho(3)=5,\rho(5)=1\). Whenever \((i,j,k)\)
appears below, it is one of
\((1,3,5),(3,5,1),(5,1,3)\), so \(j=\rho(i)\) and
\(k=\rho^2(i)\). Put
\[
 r=j+1,\qquad M_i=\{i,j+1,7,9\},\qquad \ell(M_i)=2.
\]
Whenever \((p,q)\) appears it is either \((7,9)\) or \((9,7)\).
For distinct player indices, \(x(jkq)\) means
\(x_j+x_k+x_q\), and similarly for other displayed strings.
\end{hbdefinition}

\begin{hbdefinition}{def:p18-lucas-regions}{All regions and their boundaries}
For every cyclic \((i,j,k)\), define
\[
\begin{aligned}
 E_i&=\{x\in B:x_j=x_k=1,\ x_i<1,\ x_7+x_9<1\},\\
 F_{jk}&=\{x\in B:x_j=x_k=1,\ x_7+x_9\geq1\}\setminus C.
\end{aligned}
\]
For the two choices of \((p,q)\), define
\[
 F_p=\{x\in B:x_p=1,\ x_q<1,\
             x(35q)\geq2,\ x(51q)\geq2,\ x(13q)\geq2\}\setminus C.
\]
The remaining pieces are
\[
 F_{79}=\{x\in B:x_7=x_9=1\}\setminus C,\qquad
 F_{135}=\{x\in B:x_1=x_3=x_5=1\}\setminus C.
\]
Put
\[
\begin{aligned}
 E&=E_1\cup E_3\cup E_5,\\
 F&=F_{35}\cup F_{51}\cup F_{13}\cup F_7\cup F_9
                              \cup F_{79}\cup F_{135},\\
 R&=B\setminus(C\cup E\cup F).
\end{aligned}
\]
The constituent \(F\) pieces may overlap. The weak inequalities
in their definitions include their stated boundary faces.
The pair subscript is unordered: \(F_{ab}=F_{ba}\) for distinct
\(a,b\in H\).

For \(i\in H\) and \(t,u,v\in[0,1]\), let \(e_i(t,u,v)\in B\)
have odd coordinates
\[
 x_i=t,\quad x_h=1\ (h\in H\setminus\{i\}),\quad x_7=u,\quad x_9=v,
\]
and let every even coordinate be the complement of its paired
odd coordinate. This specifies all ten coordinates.
\end{hbdefinition}

\begin{hblemma}{thm:p18-lucas-region-partition}{proved=annex:p18.lucas.region.partition}{A disjoint five-region partition}
The five sets
\[
                 A\setminus B,\quad R,\quad C,\quad E,\quad F
\]
are pairwise disjoint and cover \(A\). The \(E_i\) are pairwise
disjoint, nonempty, and have the exact coordinates
\[
 E_i=\{e_i(t,u,v):0\leq t<1,\ u\geq0,\ v\geq0,\ u+v<1\}.
\]
For each \(x\in R\), at most one of \(x_1,x_3,x_5\) equals 1,
\(x(O)<4\), and at least one of \(x_7,x_9\) is less than 1.
If \(x_p=1\), then at least one of
\[
                        x(35q),\quad x(51q),\quad x(13q)
\]
is strictly less than 2.
\end{hblemma}

\begin{hblemma}{thm:p18-lucas-core-outside}{proved=annex:p18.lucas.core.outside}{The core dominates every point outside the paired cube}
For \(x\in A\setminus B\), choose an even \(h\) such that
\(x_{h-1}+x_h<1\), and put
\[
 \epsilon=\frac{1-x_{h-1}-x_h}{2},\qquad
 y=(x_{h-1}+\epsilon)c^0+(x_h+\epsilon)c^h.
\]
Then \(y\in C\) and \(y\succ_{\{h-1,h\}}x\). Consequently
\[
                         A\setminus B\subseteq\operatorname{Dom}C.
\]
\end{hblemma}

\begin{hblemma}{thm:p18-lucas-core-remainder}{proved=annex:p18.lucas.core.remainder}{The core dominates the remaining interior region}
Every \(x\in R\) is dominated by an explicitly constructible
point of \(C\). The exhaustive constructions use these cases.
If all five odd coordinates are less than 1, improve all five
to a sum of 4 and dominate via \(O\). If \(x_i=1\) and the
other four odd coordinates are less than 1, keep \(y_i=1\),
improve those four to a sum of 3, and use \(O\setminus\{i\}\).
If \(x_p=1\), select an odd triple consisting of two members
of \(H\) and \(q\), each with coordinate less than 1 and with
sum less than 2. Improve that triple to a sum of 2, set the
two remaining odd coordinates to 1, and use that triple.

Precisely, whenever a selected set \(T\) has \(m\) members,
\(x_h<1\) for every \(h\in T\), and \(s=x(T)<m-1\), the required
strict improvements are
\[
 y_h=x_h+\lambda(1-x_h),\qquad
 \lambda=\frac{m-1-s}{m-s}\in(0,1).
\]
The unselected odd coordinates are 1, and the even coordinates
are their paired complements. Thus
\[
            (A\setminus B)\cup R\subseteq\operatorname{Dom}C.
\]
\end{hblemma}

\begin{hblemma}{thm:p18-lucas-incoming-regions}{proved=annex:p18.lucas.incoming.regions}{All possible attacks into the retained regions}
If \(x\in E_i\) and \(y\in B\), then \(y\succ_S x\) can occur
only for \(S=M_i\). For each cyclic triple, a point \(y\in C\cup F\)
satisfying \(y(M_i)\leq2\) has \(y_{j+1}=0\).
It follows that
\[
                         E\cap\operatorname{Dom}(C\cup F)=\varnothing.
\]
The other exceptional region also satisfies
\[
                         F\cap\operatorname{Dom}(C\cup E\cup F)=\varnothing.
\]
These assertions include equality on the feasibility planes
and on every defining boundary of \(F\).
\end{hblemma}

\begin{hblemma}{thm:p18-lucas-cyclic-dominance}{proved=annex:p18.lucas.cyclic.dominance}{The exact directed cycle and its feasibility plane}
The only possible directions of domination inside \(E\) are
\[
                     E_5\longrightarrow E_3\longrightarrow E_1
                                                \longrightarrow E_5.
\]
For \(x\in E_i\) and \(y=e_j(t,u,v)\in E_j\),
\[
             y\succ x
       \quad\Longleftrightarrow\quad
             t\geq u+v,\quad u>x_7,\quad v>x_9.
\]
The coalition realizing this domination is \(M_i\), whose
payment at \(y\) is \(2-t+u+v\).
There is no domination within an individual \(E_i\).
If \(t<u+v\), the point \(e_j(t,u,v)\) dominates no member of
\(E\); if \(t=u+v\), it is feasible and dominates exactly the
points of \(E_i\) below it in both coordinates 7 and 9.
Even at \(t=u+v\), it dominates none when \(u=0\) or \(v=0\).
\end{hblemma}

\begin{hbdefinition}{def:p18-lucas-upper-regions}{The three determining regions above a point}
For any \(z\in E\) and \(h\in H\), define
\[
 G_h(z)=\{e_h(t,u,v):u>z_7,\ v>z_9,\ u+v\leq t<1\}.
\]
For \(X\subseteq E\), define its predecessors and successors
inside \(E\) by
\[
\begin{aligned}
 \operatorname{Pred}_E(X)&=\{y\in E:\exists x\in X\ (y\succ x)\},\\
 \operatorname{Dom}_E(X)&=E\cap\operatorname{Dom}X.
\end{aligned}
\]
These are existential set images; they impose no common
dominating witness for every member of \(X\).
\end{hbdefinition}

\begin{hblemma}{thm:p18-lucas-upper-regions}{proved=annex:p18.lucas.upper.regions}{Nonemptiness and all geometric witness identities}
For every \(z\in E\), all three \(G_h(z)\) are nonempty. For
every cyclic \((i,j,k)\),
\[
\begin{aligned}
 \operatorname{Pred}_E(G_i(z))&=G_j(z),\\
 \operatorname{Pred}_E(G_j(z))&=G_k(z),\\
 \operatorname{Pred}_E(G_k(z))&=G_i(z),\\
 G_i(z)&\subseteq\operatorname{Dom}_E(G_j(z)).
\end{aligned}
\]
If \(z\in E_j\), its exact predecessor set is \(G_k(z)\).
The witnesses for these statements can always be chosen with
their active coordinate equal to the sum of coordinates 7 and
9. No supremum, closedness of a proposed solution, or limiting
point at \(u+v=1\) is required.
\end{hblemma}

\begin{hbtheorem}{thm:p18-lucas-local-nonexistence}{proved=annex:p18.lucas.local.nonexistence}{No stable solution exists even relative to \(E\)}
There is no \(K'\subseteq E\) that is internally stable and
dominates every member of \(E\setminus K'\). In particular,
allowing \(K'\) to be infinite, nonclosed, or nonmeasurable
does not repair the obstruction.
\end{hbtheorem}

\begin{hbtheorem}{thm:p18-lucas-no-solution}{proved=annex:p18.lucas.no.solution}{The exact ten-player nonexistence theorem}
Lucas's game \((N,\ell)\) has no NM solution. Its domination
argumentation framework, whose arguments are all imputations
in \(A_\ell\), therefore has no stable extension.
The nonempty core
\[
 C=\operatorname{conv}\{c^0,c^2,c^4,c^6,c^8,c^{10}\}
\]
is unchanged. Thus existence of undominated imputations does
not imply existence of an internally and externally stable
set of imputations.
\end{hbtheorem}

\begin{hbtheorem}{thm:p18-lucas-positive-cover}{proved=annex:p18.lucas.positive.cover}{The superadditive game has the same nonexistence}
For the partition cover \(\widehat\ell\) of
\cref{def:p18-superadditive-cover},
\[
          \widehat\ell(S)=\ell(S)\quad\hbox{whenever }\ell(S)>0.
\]
In particular \(\widehat\ell(N)=5\), and singleton values remain
zero. The cover is superadditive and has exactly the same
imputation set, core and complete domination graph as \(\ell\).
It therefore has no NM solution. This is a ten-player
counterexample within the superadditive domain used for
cooperative games in \cref{thm:p18-nm-argumentation}.
\end{hbtheorem}

\paragraph{Worked boundary calculations.}
Consider the direction \(E_3\longrightarrow E_1\). Let
\[
 x=e_1(0,\tfrac18,\tfrac18),\qquad
 y=e_3(\tfrac12,\tfrac14,\tfrac14).
\]
The target has coordinates \(x_1=x_4=0\); the attacker has
\(y_1=1\), \(y_4=\tfrac12\), and strictly larger coordinates
7 and 9. Its coalition \(\{1,4,7,9\}\) has sum
\(1+\tfrac12+\tfrac14+\tfrac14=2\), exactly its value.
Thus \(y\succ x\) on the weak feasibility boundary. Replacing
\(y_3=\tfrac12\) by \(\tfrac12-\eta\), for
\(0<\eta<\tfrac12\), raises this payment to \(2+\eta\) and
removes every outgoing attack by that point inside \(E\).
The point \(e_3(\tfrac12,0,\tfrac14)\) is feasible but has no
outgoing attack inside \(E\), since no nonnegative coordinate
7 can be strictly improved by zero.

For \(z=x\), choosing \(u=v=\tfrac5{16}\) and
\(t=\tfrac58\) gives a point in each \(G_h(z)\).
At the other boundary, \(e_1(\tfrac12,\tfrac12,\tfrac12)\)
belongs to \(F_{35}\), since its odd sum is \(\tfrac72<4\)
and its last two odd coordinates sum to 1. It is not in \(E_1\).
The point \(e_1(1,\tfrac14,\tfrac14)\) belongs to \(F_{135}\);
the point \(e_1(1,\tfrac12,\tfrac12)\) is in \(C\).
These examples show why replacing any of the stated strict
or weak boundary conditions changes the partition used in
the proof.
